\section{Tres técnicas principales de detección de vulnerabilidades}

\subsection{SAST: prueba estática de seguridad de aplicaciones}

\begin{itemize}
    \item Técnica de prueba de \textbf{caja blanca}, analiza el código fuente y los binarios
    \item Se ejecuta en las primeras etapas del SDLC, cuando el costo de detección y corrección es más bajo
    \item Puede identificar: SQL injection, command injection, XSS
    \item Técnica central: exploración de la base de código basada en coincidencia de patrones
\end{itemize}

\begin{knowledgebox}{Concepto de seguridad ``Shift Left''}
SAST plasma el concepto de seguridad ``shift left'', que consiste en adelantar las revisiones de seguridad lo más posible hacia las primeras etapas del ciclo de desarrollo. Detectar la vulnerabilidad en la etapa de escritura del código tiene un costo de corrección mucho menor que detectarla después, ya en producción. La IA hace que este concepto sea más fácil de poner en práctica que antes.
\end{knowledgebox}

\subsection{DAST: prueba dinámica de seguridad de aplicaciones}

\begin{itemize}
    \item Técnica de prueba de \textbf{caja negra}, simula el comportamiento de un atacante real
    \item Puede aplicarse a lo largo de todo el SDLC, con una tasa de falsos positivos más baja
    \item Puede identificar: SQL injection, broken authentication, XSS
    \item Técnica central:
    \begin{itemize}
        \item Input fuzzing (prueba de fuzzing de entradas)
        \item Session token manipulation (manipulación de tokens de sesión)
        \item Configuration/header testing (prueba de configuración/encabezados)
        \item Brute force rate-limit tests (pruebas de límite de tasa contra fuerza bruta)
    \end{itemize}
\end{itemize}

\subsection{SCA: análisis de composición de software}

\begin{itemize}
    \item Analiza a fondo los paquetes de código abierto (OSS) que usa la aplicación
    \item Analiza los administradores de paquetes, infrastructure-as-code y las imágenes descargadas
    \item Técnica central:
    \begin{itemize}
        \item Analizar las dependencias a partir de los metadatos de los paquetes
        \item Resolución de dependencias transitivas (transitive dependency resolution)
        \item Comparación con bases de datos de vulnerabilidades conocidas
        \item Exploración de binarios/artifacts
    \end{itemize}
\end{itemize}

\subsection{Resumen de la sección}

Las tres técnicas SAST, DAST y SCA se complementan entre sí: SAST detecta vulnerabilidades en etapas tempranas mediante el análisis del código, DAST simula ataques en tiempo de ejecución, y SCA analiza la seguridad de las dependencias de terceros. Una estrategia de seguridad completa requiere combinar las tres.
